\section{Artifact and Reproduction Guide}
\label{app:artifact}
An artifact should connect each reported result to an executable command,
configuration, and raw output. Record immutable source revisions and
dependency versions. Document the time and resources needed for a small
sanity check and for the full evaluation.

\subsection{A minimal execution record}
\begin{lstlisting}[language=STARPython,caption={Illustrative record format;
replace fields with the study's actual schema.},label={lst:record}]
record = {
    "subject": "replace-with-id",
    "revision": "replace-with-commit",
    "seed": 0,
    "status": "ok",
    "raw_output": "results/run.json",
}
\end{lstlisting}

Use separate statuses for timeouts, crashes, invalid inputs, and completed
runs. Preserve enough evidence to distinguish a harness failure from a
failure of the system being studied.

\subsection{Release manifest}
\begin{itemize}
  \item Source and configuration with explicit licenses.
  \item Data provenance, access rules, and preprocessing steps.
  \item Hardware and software environment specifications.
  \item Commands that regenerate each table and figure.
  \item Raw outputs, labeling records, and known limitations.
\end{itemize}

% In this short demonstration appendix, begin the independent editing guide
% at the top of the second column. Remove this break in a longer manuscript.
\newpage
\section{Working with the STAR Template}
Edit \code{paper-info-double.tex} for paper metadata and the files in
\code{sections/double/} for content. Store reusable notation in
\code{macros.tex}. The \code{starpaper.cls} file controls both layouts;
the two-column option switches the typeface, geometry, citation style,
and title treatment together.

Use \verb|figure| for a column-width illustration and
\verb|figure*| for an illustration spanning both columns. Full-width
floats normally appear at the top of a later page. Refer to every figure
in the text, and keep its labels readable at its final printed size.

\subsection{Anonymous and grayscale modes}
The \code{anonymous} class option removes the lab logo, author block,
affiliation, emails, identifying running furniture, acknowledgments, and
author metadata managed by the class. It does not inspect prose or
third-party assets. Review links, self-references, and supplementary files
before an anonymous submission.

The \code{monochrome} option changes the class colors and built-in logo to
grayscale. External figures retain their original colors and must be
checked separately.

\subsection{Using an official venue style}
Use this original lab style for manuscripts and preprints. Use the official
venue template for conference submissions. Reuse section files and the
bibliography, load the required packages, and replace lab-specific commands.
Remove identifying marks from anonymous manuscripts.
